\section{Métricas experimentales clave}
\subsection{Observaciones de Iteration y win rate}
Weston divide el entrenamiento de self-rewarding en varias iteraciones; cada ronda combina tareas de Self-Instruct, verified response y judge. Los experimentos muestran que la primera ronda de entrenamiento ya le permite al modelo obtener un win rate considerable en Alpaca Eval 2, y que al agregar el meta judge el win rate mejora aún más, evitando el plateau que originalmente aparecía en la tercera ronda (SRT 1908-2027, 2784-2790).

\begin{table}[H]
\centering
\begin{tabular}{p{4cm}p{4cm}p{7cm}}
\toprule
Etapa de iteración & Estimación del tamaño de muestra & Observación \\
\midrule
Seed Iteration & Pocos verified prompts & El modelo domina rápido el tipo de tarea, pero la calidad del reasoning aún depende del judge. \\
Self-Rewarding Iteration 1 & Self-Instruct + verified response & El win rate mejora inicialmente en Alpaca Eval 2; el modelo aprende a autoevaluarse. \\
Self-Rewarding Iteration 2 & Tareas nuevas + cross-check & Se verifican los errores de manera repetida y el reward se desplaza hacia la dirección correcta. \\
Meta-Judge Iteration & Se agrega el judge de SelfT evals & Se resuelve el plateau, la confiabilidad del verified response se vuelve cuantificable y el win rate sube de manera constante. \\
\bottomrule
\end{tabular}
\caption{Observaciones de iteration y win rate presentadas en la clase (qualitative).}
\end{table}

\subsection{SelfT evaluators y otras fuentes de datos}
Además de Alpaca Eval 2, Weston también menciona SelfT, Self-Feeding Chatbot y generic instruction data (matemáticas/código/juicio/conversación). Estas fuentes de datos se complementan entre sí: SelfT aporta judge supervision, Self-Feeding aporta conversational feedback y Self-Instruct aporta new tasks, y forman los cuatro puntos de apoyo necesarios para el «autoentrenamiento del modelo».

\begin{importantbox}{Los cuatro puntos de las fuentes de datos}
\begin{itemize}
    \item Self-Instruct: genera instruction → input → output.
    \item Self-Feeding Chatbot: genera verified response mediante la conversación.
    \item SelfT evaluators: aporta etiquetas simuladas de judge.
    \item Alpaca Eval 2: observación del win rate en el benchmark, que se usa para guiar el modelo de reward.
\end{itemize}
\end{importantbox}

\subsection{Resumen de la sección}
Las distintas iteraciones intensifican progresivamente tanto las fuentes de datos como el judge: Self-Instruct se encarga de las tareas, Self-Feeding aporta la verificación, SelfT entrena al judge y Alpaca Eval 2 aporta el win rate final; los cuatro elementos trabajan juntos para mejorar de manera constante el desempeño en reasoning.

